\documentclass[11pt]{article}

\usepackage[margin=1in]{geometry}
\usepackage{amsmath,amssymb,amsthm}
\usepackage{mathtools}
\usepackage{enumitem}
\usepackage{booktabs}
\usepackage{array}
\usepackage{float}
\usepackage{hyperref}
\hypersetup{colorlinks=true,linkcolor=blue,citecolor=blue,urlcolor=blue,
  pdftitle={Exact count-n EFX instances for additive valuations beyond m=n+2 and a bounded counting verification},
  pdfsubject={EFX allocations, counting, computer-assisted verification}}
\setlength{\emergencystretch}{3em}

\theoremstyle{plain}
\newtheorem{theorem}{Theorem}[section]
\newtheorem{lemma}[theorem]{Lemma}
\newtheorem{proposition}[theorem]{Proposition}
\newtheorem{corollary}[theorem]{Corollary}
\newtheorem{conjecture}[theorem]{Conjecture}

\theoremstyle{definition}
\newtheorem{definition}[theorem]{Definition}
\newtheorem{remark}[theorem]{Remark}
\newtheorem{cverified}[theorem]{Computer-verified result}

\newcommand{\efxo}{\mathrm{EFX}_0}
\newcommand{\Z}{\mathbb{Z}}

\title{Exact count-\texorpdfstring{$n$}{n} EFX instances for additive valuations beyond
\texorpdfstring{$m=n+2$}{m=n+2}\\and a bounded counting verification}
\author{% TODO-AUTHOR: fill in author name(s), affiliation, and email before submission
}
\date{}

\begin{document}
\maketitle

\begin{abstract}
Neoh and Teh initiated the systematic study of how many allocations satisfy
envy-freeness up to any good (EFX).  Their Conjecture~6.1 predicted at least
$n$ EFX allocations for general monotone valuations, but the 2026 monotone
submodular counterexamples of Akrami et al.\ refuted that statement.  Its additive
specialization remains open.  We give positive additive instances having
exactly $n$ labeled EFX allocations at $(n,m)=(4,7),(5,8),(6,9)$ and at
$(4,8),(4,9),(4,10)$.  Thus the minimum over additive instances is at most
$n$ at these six cells; equality remains conditional on the additive lower-bound
conjecture.  The allocations in these instances
are explained by a generalized parking criterion: for an EFX allocation of all
but one good, an explicit observer--host inequality is necessary and sufficient
for placing the remaining good at a chosen host.  Finally, a slot-protected
counterexample-guided computation reports that every positive integer $(4,7)$
instance with values in $\{1,\ldots,4\}$ has at least four EFX allocations.  We
prove the soundness of the slot encoding and provide exact instance verifiers,
parking-slack certificates, and a reconstruction and replay package.  Because
the historical run did not retain its final checkpoint, we regenerated an
alternative complete terminal pool; the fresh CP-SAT master is infeasible at
pool size $1{,}240$, and an independently encoded CNF master is also UNSAT
under CaDiCaL.  The artifact archives CaDiCaL's binary DRAT proof,
independently verified by \texttt{drat-trim} \cite{DRATtrim}.  The bounded result depends on
the stated CNF translation and proof checker; it is not an unrestricted
counting theorem.
\end{abstract}

% ===========================================================================
\section{Introduction}
% ===========================================================================

Let $N=[n]$ be a set of agents and $M=[m]$ a set of indivisible goods.  Agent
$i$ has an additive valuation $v_i\colon 2^M\to\mathbb{R}_{\ge0}$, with
$v_i(S)=\sum_{g\in S}v_i(g)$.  An allocation is a labeled partition
$X=(X_1,\ldots,X_n)$ of $M$.

\begin{definition}\label{def:efx}
An allocation $X$ is \emph{$\efxo$} if, for every pair $i\ne j$ and every
$g\in X_j$,
\[
  v_i(X_i)\ \ge\ v_i(X_j\setminus\{g\}).
\]
\end{definition}

The subscript distinguishes the strong convention in which zero-valued goods are
also quantified over.  Every computational instance and bounded result in this
paper uses positive values, so the distinction disappears there; the parking
theorem is stated directly for $\efxo$ and permits nonnegative valuations.  We
count labeled allocations: permuting bundles among agents generally produces a
different allocation.

Suksompong \cite{Suksompong2020} obtained tight counting results for two agents.
Neoh and Teh \cite{NeohTeh2025} subsequently studied the minimum count as a
function of $(n,m)$ for general monotone valuations, using additive examples
for their upper bounds.  Their Table~1 gives the exact value $n$ at $m=n+1$,
bounds $[n,n^2]$ at $m=n+2$, and groups $m\ge n+3$ under the upper bound
$n!$ with no positive lower bound uniform across that entire regime.  They
conjectured at least $n$ allocations without an additivity restriction.
Akrami et al.\ \cite{Akrami2026SAT} have since
constructed monotone submodular instances with no EFX allocation for
$n\ge3$ and $m\ge n+5$, refuting that original conjecture.

\begin{conjecture}[additive specialization of Neoh--Teh
\cite{NeohTeh2025}, Conjecture~6.1]\label{conj:nt}
With $n$ agents and additive valuations, there are at least $n$ $\efxo$
allocations.
\end{conjecture}

Conjecture~\ref{conj:nt} is the still-open additive restriction, not the
now-false general-monotone statement.  It includes additive EFX existence as a
special consequence.  For clarity, call an additive instance \emph{count-$n$}
if its complete labeled allocation set contains exactly $n$ $\efxo$
allocations.  We do not call such an instance tight unless a matching additive
lower bound is available.

\paragraph{Contributions.}
The paper has three claims.
\begin{enumerate}[label=(\roman*),leftmargin=2.3em]
\item We give explicit count-$n$ instances at six cells:
\[
 (4,7),\ (5,8),\ (6,9),\ (4,8),\ (4,9),\ (4,10).
\]
Two independently written exact verifiers and a third vectorized implementation
enumerate all $n^m$ labeled allocations and agree on every count.  These are
upper bounds on the additive minimum count; they do not prove
Conjecture~\ref{conj:nt}.
\item We prove a generalized parking theorem.  It gives an exact
observer--host inequality for when an unallocated good may be added to a bundle
while preserving $\efxo$.  Machine-checkable slack matrices show that all $n$
parkings in each displayed instance satisfy the theorem.  Universal-minimum
parking is only a sufficient special case, and is not assumed for the family.
\item We formalize a slot-protected counting-CEGAR master and prove its soundness.
CP-SAT and an independently encoded threshold-CNF master both report UNSAT
over the finite box $\{1,\ldots,4\}^{4\times7}$ for the hypothesis ``at most
three $\efxo$ allocations.''  The artifact audits the historical checkpoint,
supplies a complete alternative terminal pool generated by atomic
reconstruction, and archives both masters, their replay programs, and an
independently checked binary DRAT certificate for the CNF master.
\end{enumerate}

The instance upper bounds and bounded lower bound are logically separate.  In
particular, the printed $(4,7)$ instance has values well outside
$\{1,\ldots,4\}$, so the two computations cannot be combined into an
unrestricted equality statement.

\paragraph{Related computational work.}
For additive valuations, EFX exists for at most three agents
\cite{ChaudhuryThree}.  Mahara \cite{Mahara2107} proves EFX existence when
$m\le n+3$, so existence is unconditional in our three diagonal cells even for
general monotone valuations.  Together these results make $(4,8)$ the first
additive parameter cell not covered by a general existence theorem.
Akrami et al.\ \cite{Akrami2026SAT} refute general-monotone EFX existence from
$m=n+5$ onward, while explicitly leaving additive valuations outside their
counterexamples.  Their work also gives a recent computer-assisted precedent:
a Lean-checked theoretical foundation for the SAT encoding, two SAT solvers,
and a DRAT-checked UNSAT proof.  Our finite result concerns bounded additive
integer valuations and now also retains an independently checked DRAT proof,
but no proof-assistant formalization of its mathematical encoding.  The
witness search is also distinct
from Br\^anzei's simulated-annealing search for an EFX allocation of a fixed
instance \cite{Branzei2025}: here the search varied valuation matrices to minimize
the exact allocation count.

% ===========================================================================
\section{Six count-\texorpdfstring{$n$}{n} cells}\label{sec:instances}
% ===========================================================================

For an additive instance $I$, write $\#\efxo(I)$ for its number of complete
labeled $\efxo$ allocations, and let the \emph{additive minimum count} be the
minimum of this quantity over additive instances with the same $(n,m)$.  An
explicit instance with $\#\efxo(I)=n$ proves that this additive minimum is at
most $n$, regardless of whether the additive lower bound is $0$, lies in
$\{1,\ldots,n-1\}$, or equals $n$.

\begin{cverified}[count-$n$ instances]\label{cv:instances}
There are positive additive instances with exactly $n$ $\efxo$ allocations at
$(4,7)$, $(5,8)$, and $(6,9)$, and positive additive instances with exactly four
$\efxo$ allocations at $(4,8)$, $(4,9)$, and $(4,10)$.
\end{cverified}

Rows below are agents and columns are goods.  One witness for each cell is:
\begin{align*}
(4,7):\quad & \begin{bmatrix}
30&66&69& 9&21&19&63\\
12&56&58& 9&10&24&52\\
12&33&45&20&24& 8&64\\
11&60& 1&18& 2& 6&41
\end{bmatrix}
\\[1ex]
(5,8):\quad & \begin{bmatrix}
32& 4&10&60& 8& 3&16&11\\
42& 3&57&50&51& 1& 5&10\\
54&12& 4&27& 3&20&18&11\\
33& 9&31&42&45&10&13& 5\\
45&13&52&47&49&21& 2&21
\end{bmatrix}
\\[1ex]
(6,9):\quad & \begin{bmatrix}
15&86&103&193&191&192&195&196&1\\
 6& 3& 12&105& 23& 65&  4&  2&1\\
11&68& 55& 10&152&  8&  7&  2&1\\
49& 3& 11&  2& 64&104&  8&  7&1\\
 2& 6& 31& 79& 40& 61&104&  7&1\\
53&45&  3&105&102&103&151&152&2
\end{bmatrix}.
\end{align*}

For the three remaining cells:
\begin{align*}
(4,8):\quad &
\begin{bmatrix}
11& 8&20&1& 4&1&16&1\\
16&49&32&2&16&2&60&1\\
28&28&59&1&60&7&46&8\\
10&12&24&1&59&1&17&1
\end{bmatrix}
\\[1ex]
(4,9):\quad & \begin{bmatrix}
1&57&58& 6&60& 7&19&14&15\\
1&57&11&12& 4&12& 1&21&10\\
4&59&60& 2&23& 3&28& 8&12\\
8&48&57& 1&58&34& 1& 1& 3
\end{bmatrix}
\\[1ex]
(4,10):\quad & \begin{bmatrix}
17&57&58&3&47& 1& 6&23& 3& 1\\
 2&51&57&5&58&11&12& 1&10&10\\
 1&75&15&8&25&13&19& 6&12&16\\
10&58&59&8&60& 2& 8& 7&22& 6
\end{bmatrix}.
\end{align*}

\paragraph{Verification.}
Verifier~1 enumerates assignments by a base-$n$ counter and checks bundle
bitmasks.  Verifier~2 uses a recursive good-by-good enumeration and separately
constructs observer--bundle totals and minima.  A vectorized counter supplies a
third implementation.  All arithmetic is integral.  For each observer--target
pair, the strongest $\efxo$ constraint is obtained by removing an
observer-minimum good from the target bundle, so the optimized check is
\[
  v_i(X_i)\ \ge\
  v_i(X_j)-\min_{g\in X_j}v_i(g).
\]
This optimized constraint is imposed only when $X_j$ is nonempty; an empty
target bundle has no removable-good constraint.
The second verifier is also tested against a definition-faithful loop that checks
every removable good separately.  The artifact contains the instance files,
both exact verifiers, expected outputs, and hashes.

The instances were found by directed count-minimizing search.  Search restrictions
do not enter Result~\ref{cv:instances}: once a matrix is printed, exhaustive
enumeration alone establishes its count.

\begin{table}[H]
\centering
\begin{tabular}{@{}lrrrl@{}}
\toprule
$(n,m)$ & allocations checked & $\#\efxo$ & prior additive UB & artifact instance \\
\midrule
$(4,7)$  & $4^7=16{,}384$       & $4$ & $24$  & \texttt{t47.json}\\
$(5,8)$  & $5^8=390{,}625$      & $5$ & $120$ & \texttt{ch5\_A.json}\\
$(6,9)$  & $6^9=10{,}077{,}696$ & $6$ & $720$ & \texttt{tight69\_witness.json}\\
$(4,8)$  & $4^8=65{,}536$       & $4$ & $24$  & \texttt{tb4.json}\\
$(4,9)$  & $4^9=262{,}144$      & $4$ & $24$  & \texttt{tight49\_witness.json}\\
$(4,10)$ & $4^{10}=1{,}048{,}576$ & $4$ & $24$ & \texttt{tight410\_witness.json}\\
\bottomrule
\end{tabular}
\caption{Full labeled-allocation enumerations for the six cells.  ``Prior
additive UB'' is the $n!$ upper bound witnessed by the additive examples of
Neoh and Teh \cite{NeohTeh2025}.}
\label{tab:counts}
\end{table}

% ===========================================================================
\section{A generalized parking theorem}\label{sec:parking}
% ===========================================================================

Let $f$ be a good and let $B=(B_1,\ldots,B_n)$ allocate $M\setminus\{f\}$.
For host $h$, write $B^{+h}$ for the complete allocation obtained by adding $f$
to $B_h$.  For observer $k\ne h$, define the \emph{parking slack}
\begin{equation}\label{eq:slack}
\sigma_{k,h}(B,f)
 =
 v_k(B_k)-
 \left(
   v_k(B_h)+v_k(f)
   -\min_{g\in B_h\cup\{f\}}v_k(g)
 \right).
\end{equation}
The minimum is always defined because the enlarged host bundle contains $f$.

\begin{theorem}[generalized parking criterion]\label{thm:parking}
Suppose $B$ is $\efxo$.  For a fixed host $h$, the completion $B^{+h}$ is
$\efxo$ if and only if
\[
  \sigma_{k,h}(B,f)\ge0
  \qquad\text{for every }k\ne h.
\]
Consequently all $n$ parkings are $\efxo$ if and only if every off-diagonal
entry of the parking-slack matrix is nonnegative.
\end{theorem}

\begin{proof}
Fix a host $h$.  Consider an ordered observer--target pair after adding $f$.
If the target is not $h$ and the observer is not $h$, both bundles are unchanged
from $B$, so the $\efxo$ inequalities follow from the hypothesis.  If the
observer is $h$ and the target is some $j\ne h$, the observer's own value has
weakly increased while the target is unchanged, so these inequalities also
remain valid.

It remains to consider observer $k\ne h$ and target $h$.  The observer's own
bundle is $B_k$, while the target is $B_h\cup\{f\}$.  Among all possible removed
goods, the largest remaining target value is obtained by removing a
$k$-minimum good.  Thus the complete family of $\efxo$ inequalities for this
pair is equivalent to
\[
v_k(B_k)\ge
v_k(B_h)+v_k(f)-\min_{g\in B_h\cup\{f\}}v_k(g),
\]
which is exactly $\sigma_{k,h}(B,f)\ge0$.  This proves both necessity and
sufficiency.
\end{proof}

When $B_h$ is nonempty, the condition can equivalently be written
\[
v_k(B_k)\ge
v_k(B_h)+
\max\!\left\{0,\,
v_k(f)-\min_{g\in B_h}v_k(g)\right\}.
\]
This form makes clear that a floater need not be a minimum-valued good: the base
may instead provide enough slack to absorb its excess value.

\begin{corollary}[universal-minimum special case]\label{cor:universal}
Suppose $B$ is envy-free and $v_k(f)\le v_k(g)$ for every agent $k$ and every
$g\in M\setminus\{f\}$.  Then all $n$ parkings of $f$ are $\efxo$.
\end{corollary}

\begin{proof}
Envy-freeness implies $\efxo$ under nonnegative valuations, so $B$ meets the
hypothesis of Theorem~\ref{thm:parking}.
For nonempty $B_h$, universal minimality makes the parenthesized target
remainder in \eqref{eq:slack} equal to $v_k(B_h)$; for empty $B_h$ it equals
$0$.  Envy-freeness gives the required inequality in either case, and
Theorem~\ref{thm:parking} applies.
\end{proof}

Neoh and Teh use the zero-valued-good special case of
Corollary~\ref{cor:universal} in their counting analysis
\cite{NeohTeh2025}.  The corollary is also related to, but is not being
attributed as a lemma of, the envy-graph and charity techniques in
\cite{Lipton2004,CGM2019,Berger2022}.  Theorem~\ref{thm:parking} is the direct
exact criterion needed here.

\paragraph{Application to the witnesses.}
For each matrix in Section~\ref{sec:instances}, the artifact identifies a good
$f$ and an $\efxo$ base $B$ such that every $\efxo$ allocation of that instance is one
of the $n$ parkings $B^{+h}$.  It records every integer
$\sigma_{k,h}(B,f)$ and checks it directly from the matrix.  Only the $(6,9)$
floater is a universal minimum; no such claim is needed.  Theorem
\ref{thm:parking} proves that the $n$ listed parkings are $\efxo$, while the
exhaustive verifiers establish that there are no additional allocations.
For completeness, the certified bases below are written in agent order with
one-based goods:
\[
\begin{array}{c@{\qquad}l}
(4,7)  & (\{1,5,6\},\{3\},\{7\},\{2\}),\\
(5,8)  & (\{4\},\{3\},\{1\},\{5\},\{2,6,8\}),\\
(6,9)  & (\{1,2,3\},\{4\},\{5\},\{6\},\{7\},\{8\}),\\
(4,8)  & (\{3\},\{1,2\},\{6,7,8\},\{5\}),\\
(4,9)  & (\{4,6,7,8,9\},\{2\},\{3\},\{5\}),\\
(4,10) & (\{3\},\{5\},\{2\},\{1,4,7,8,9,10\}).
\end{array}
\]

\begin{table}[t]
\centering
\begin{tabular}{@{}lcccccc@{}}
\toprule
cell & floater & base shape & min.\ EF & min.\ parking & universal min. & count\\
\midrule
$(4,7)$  & $4$ & $3111$   & $1$ & $1$ & no  & $4$\\
$(5,8)$  & $7$ & $31111$  & $3$ & $3$ & no  & $5$\\
$(6,9)$  & $9$ & $311111$ & $1$ & $1$ & yes & $6$\\
$(4,8)$  & $4$ & $3211$   & $1$ & $1$ & no  & $4$\\
$(4,9)$  & $1$ & $5111$   & $1$ & $1$ & no  & $4$\\
$(4,10)$ & $6$ & $6111$   & $1$ & $1$ & no  & $4$\\
\bottomrule
\end{tabular}
\caption{Parking certificates for the six displayed matrices.  Goods are
one-based.  The artifact records the full base bundles and observer--host slack
matrices; ``min.\ EF'' is the minimum ordinary envy-free base margin.}
\label{tab:parking}
\end{table}

% ===========================================================================
\section{Bounded counting by protected slots}\label{sec:cegar}
% ===========================================================================

We next consider the finite domain
\[
 \mathcal V_K=\{1,\ldots,K\}^{n\times m}.
\]
Let $\mathcal A=[n]^m$ be the set of labeled allocations, represented by their
good-to-agent assignment vectors, and let $E(V,A)$ denote that $A$ is $\efxo$
under valuation matrix $V$.

\subsection{The master formula}

Fix an integer target lower bound $2\le c\le|\mathcal A|$ and a finite pool
$P\subseteq\mathcal A$.  Introduce
$c-1$ slot allocations $S_1,\ldots,S_{c-1}\in\mathcal A$.  The protected-slot
master is
\begin{equation}\label{eq:master}
\begin{split}
 V\in\mathcal V_K,\qquad C(V),\qquad
 \bigwedge_{A\in P}
 \left(
   \neg E(V,A)\ \lor\
   \bigvee_{t=1}^{c-1}(A=S_t)
 \right),
\end{split}
\end{equation}
where $C(V)$ is an orbit-complete symmetry-breaking condition.  Our CP-SAT
master uses row/column DoubleLex \cite{Flener2002}; relabeling agents or goods
preserves the number of $\efxo$ allocations.
DoubleLex is orbit-complete: choose a lexicographically maximal row-major
matrix in any row/column-permutation orbit.  An out-of-order adjacent row or
column could be swapped to make that matrix larger, so its rows and columns
must both be lexicographically ordered.

For integer valuations, non-EFX of a fixed allocation $A$ is encoded exactly.
Writing $A_g=j$ when good $g$ is assigned to $j$, we have
\begin{equation}\label{eq:kill}
\neg E(V,A)
\iff
\bigvee_{\substack{g\in M,\ j=A_g\\i\ne j}}
\left[
 \sum_{q:A_q=i}v_{i,q}+v_{i,g}+1
 \le
 \sum_{q:A_q=j}v_{i,q}
\right].
\end{equation}
The $+1$ is the integral form of strict strong envy.  Equality $A=S_t$ is the
conjunction of the $m$ corresponding assignment equalities.

\begin{lemma}[protected-slot soundness]\label{lem:slots}
Assume $C$ contains at least one representative of every row/good relabeling
orbit.  If the master \eqref{eq:master} is infeasible for any pool $P$, then
every $V\in\mathcal V_K$ has at least $c$ $\efxo$ allocations.
\end{lemma}

\begin{proof}
Suppose instead that some $V\in\mathcal V_K$ has $r<c$ $\efxo$ allocations.
Relabel its rows and columns to obtain a representative satisfying $C$; the
count remains $r$.  Assign its $r$ $\efxo$ allocations to $r$ protected slots
and fill any unused slots arbitrarily.  For every $A\in P$, either $A$ is not
$\efxo$, satisfying the first disjunct, or it is one of the at most $r$
protected allocations, satisfying a slot disjunct.  Hence the master is
feasible, a contradiction.
\end{proof}

The slot assignments need not be distinct or themselves constrained to be
$\efxo$.  Allowing arbitrary slots only relaxes the master, so it cannot create
a false infeasibility conclusion.

\paragraph{Independent cardinality form.}
Because $P$ is a set of distinct allocations, existentially eliminating the
slots in \eqref{eq:master} gives the equivalent constraint
\[
  \sum_{A\in P}\mathbf 1[E(V,A)]\le c-1.
\]
Indeed, slots can cover every surviving pooled allocation exactly when there
are at most $c-1$ distinct survivors.  The independent Z3 re-encoding uses
this cardinality form and one-hot finite-domain value variables, rather than
sharing the CP-SAT slot-model builder.

\paragraph{Independent CNF form.}
For $K=4$, a second checker retains the literal slots and independently
translates the master to CNF.  It represents each value by threshold bits
$y_{i,q,t}=\mathbf 1[v_i(q)\ge t]$ for $t=2,3,4$, with
$y_{i,q,4}\Rightarrow y_{i,q,3}\Rightarrow y_{i,q,2}$ and
$v_i(q)=1+\sum_{t=2}^4y_{i,q,t}$.  Fix a disjunct of
\eqref{eq:kill}, let $a=|A_i|$, and let
$b=|A_j\setminus\{g\}|$.  Its strict-envy inequality is equivalent to
\begin{equation}\label{eq:threshold-kill}
 \sum_{\substack{q\in A_j\setminus\{g\}\\t=2,3,4}}y_{i,q,t}
 +\sum_{\substack{q\in A_i\\t=2,3,4}}(1-y_{i,q,t})
 \ \ge\ 4a-b+1.
\end{equation}
Indeed, the left side is
$v_i(A_j\setminus\{g\})-b+4a-v_i(A_i)$.
The checker guards PySAT's satisfiability-preserving sequential-counter
encoding \cite{PySAT,Sinz2005} of \eqref{eq:threshold-kill} by the corresponding
violation literal, gives every good in each slot a one-hot owner, and encodes
the first strict coordinate of every DoubleLex comparison explicitly.  Its
DIMACS formula therefore shares the terminal pool and mathematical
specification with the CP-SAT and Z3 checks, but not either model builder.

\subsection{CEGAR loop and bounded result}

Starting from a pool, CP-SAT solves \eqref{eq:master}.  If the master is
feasible, an exact slave enumerates the candidate valuation's $\efxo$
allocations.  A candidate with fewer than $c$ is a counterexample.  Otherwise
the slave adds allocations to the pool.  Since a feasible master permits at
most $c-1$ pooled allocations to remain $\efxo$, returning at least $c$
candidate allocations necessarily supplies a new pool member.  Thus the pool
strictly grows inside the finite set $\mathcal A$, so only finitely many such
refinements are possible.  If the master
becomes infeasible, Lemma~\ref{lem:slots} yields the bounded lower bound.

\begin{cverified}[bounded $(4,7)$ count]\label{cv:bounded}
A CP-SAT \cite{ORtools} execution of the protected-slot master with $n=4$,
$m=7$, $K=4$, and $c=4$ reports infeasible.  A fresh reconstruction terminates
with a complete pool of $1{,}240$ distinct allocations.  The independently
encoded literal-slot CNF for that pool is also UNSAT under CaDiCaL
\cite{CaDiCaL}.  CaDiCaL's archived binary DRAT proof is accepted by the
independent \texttt{drat-trim} checker.  Conditional on the reported CNF
translation and proof checker, every positive integer $(4,7)$ instance with
entries in $\{1,\ldots,4\}$ has at least four $\efxo$ allocations.
\end{cverified}

The implementation creates three slots.  For each pooled allocation it adds
the exact disjunction \eqref{eq:kill} or equality with one of those slots.  It
uses no valuation-type restriction.  A separate fixed-instance test checks the
off-by-one boundary: protecting three distinct $\efxo$ allocations is feasible,
whereas pooling all four allocations of the fixed count-$4$ witness is
infeasible.

\subsection{Reproducibility and trust boundary}\label{sec:artifact}

The accompanying supplementary artifact freezes the claim-bearing inputs and
source, with pinned external solver and proof-checker revisions.  Its entry points are
\path{repro/counting_witnesses},
\path{artifacts/generalized_parking}, and
\path{repro/countcegar_k4}.  It includes:
\begin{itemize}[leftmargin=1.6em]
\item the six valuation matrices, expected counts, two pure-Python exact
verifiers, a vectorized exact counter, and a definition-faithful regression;
\item the generalized-parking bases, floaters, complete slack matrices, and a
standard-library certificate checker;
\item an audit of the historical checkpoint, an atomic reconstruction program,
the complete $1{,}240$-allocation terminal pool, run metadata, logs, and a
SHA-256 manifest; and
\item an independently written Z3 \cite{Z3} master checker and SMT-LIB2 exporter
for any reconstructed terminal pool, together with the threshold-CNF checker,
differential tests, frozen DIMACS input, compressed binary DRAT proof, pinned
tool revisions, and proof-verification wrapper described above.
\end{itemize}
The deterministic release workflow under \path{supplement} binds this source
and PDF to those files in a SHA-256 manifest.  It intentionally excludes
exploratory results that are not claims of this paper.

The historical result reports infeasibility at pool size $1{,}295$, but its
periodic checkpoint contains only $1{,}250$ entries, $1{,}134$ distinct.  The
last $45$ entries cannot be recovered from the historical log.  The
reconstruction starts from the distinct checkpoint entries and writes every
subsequent pool atomically.  Any terminal pool it produces therefore determines
an alternative complete terminal master for the same finite claim; it does not
reconstruct the missing historical allocations.

The fresh run begins with the $1{,}134$ distinct checkpoint allocations, makes
seven feasible master calls that add $106$ new allocations, and returns
infeasible on its eighth call.  It uses CP-SAT 9.15 with six workers and random
seed $1$ and takes $1{,}060.725$ seconds in the recorded environment.  The complete
pool is written atomically; its SHA-256 begins
\texttt{af85029ca3ec}, with the full digest recorded in the artifact manifest.
A subsequent one-call replay of the frozen pool again returns infeasible and
exports the exact CP-SAT model and response; the artifact keeps the
pool-generation and export-replay records separately.  The independent CNF has
$282{,}029$ variables and $600{,}624$ clauses.  CaDiCaL 1.9.5, called through
PySAT, returns UNSAT in $68.048$ wall-clock seconds in the recorded environment.
The archived DIMACS file has SHA-256 prefix \texttt{ccc2aa45a571}; a cold
replay that loads this frozen DIMACS without rebuilding the master reproduces
UNSAT and the same solver counters.  A direct CaDiCaL 1.9.5 invocation writes
a $134{,}426{,}255$-byte binary DRAT proof.  Official \texttt{drat-trim}
verifies the archived proof in $90.182$ seconds using $28{,}212{,}699$
resolution steps; its core uses no RAT lemmas.  The proof is stored under
\path{repro/countcegar_k4/certificates} in a $38{,}753{,}120$-byte
\texttt{xz} archive whose compressed and uncompressed SHA-256 digests are
recorded in the manifest.

The formal content of Lemma~\ref{lem:slots} and equations
\eqref{eq:master}--\eqref{eq:kill} can be checked independently of either
solver.  The universal bounded conclusion additionally trusts that the
reported masters were built as specified.  Agreement between the independently
built CP-SAT and CNF masters materially reduces the encoding risk, while the
checked DRAT trace removes reliance on CaDiCaL's raw UNSAT verdict.  The
remaining computational trust boundary is the auditable CNF translation and a
correct proof checker.  Akrami et al.\ \cite{Akrami2026SAT} additionally
formalize their theoretical SAT foundation in Lean.

% ===========================================================================
\section{Scope and limitations}
% ===========================================================================

\paragraph{Upper bounds, not tightness.}
The six explicit matrices prove upper bounds on the minimum over additive
instances.  Calling them ``tight'' would require a matching lower bound for
arbitrary additive valuations.  That lower bound is precisely the content
predicted by Conjecture~\ref{conj:nt} and is not proved here.  The original
general-monotone conjecture is already false \cite{Akrami2026SAT}.

\paragraph{Bounded means bounded.}
Result~\ref{cv:bounded} concerns positive integer entries in
$\{1,\ldots,4\}$ only.  No theorem bounds the magnitude of a minimum-count
instance, and the computation says nothing about larger or zero entries.

\paragraph{What is and is not formal.}
The generalized parking theorem and protected-slot soundness lemma are ordinary
mathematical proofs.  The fixed-instance counts are exhaustive computations
that are inexpensive to repeat with independent code.  The bounded result is a
substantially larger solver computation.  The fresh terminal pool is included
and makes its exact input replayable; the historical terminal input remains
unrecoverable.  A proof-producing solver certificate is included and
independently checked; no proof-assistant formalization is claimed.

\paragraph{Search provenance.}
Simulated annealing and constraint-guided local search were used to discover
the matrices, in the spirit of recent computational EFX search
\cite{Branzei2025}.  Discovery heuristics are not used in their verification
and are outside the paper's theorem claims.

% ===========================================================================
\section{Conclusion}
% ===========================================================================

The six count-$n$ instances show that very small EFX allocation sets already
occur at several cells beyond $m=n+2$.  Their common, rigorously stated
mechanism is not universal-minimum parking but the observer-specific slack
criterion of Theorem~\ref{thm:parking}.  The protected-slot construction gives
a general way to turn an existence CEGAR master into a bounded counting master,
and the $(4,7)$, $K=4$ computation supplies its first test case here.

The main open mathematical question addressed here is the additive
specialization: prove or refute a lower bound of $n$ for arbitrary additive
valuations.  The original general-monotone question has been answered
negatively \cite{Akrami2026SAT}.  The present results provide explicit additive
upper-bound targets, a local completion lemma, and a replayable finite
verification without claiming that they settle the additive question.

\section*{Acknowledgments}
The computational campaign used independently written enumeration, constraint,
and structure-checking programs, together with a multi-agent assistant workflow.
All mathematical and computational claims included here are scoped to the
checks described in the paper and artifact.

\begin{thebibliography}{99}

\bibitem{Suksompong2020}
W.~Suksompong.
\newblock On the number of almost envy-free allocations.
\newblock \emph{Discrete Applied Mathematics}, 284:606--610, 2020.

\bibitem{NeohTeh2025}
T.~Y. Neoh and N.~Teh.
\newblock Understanding EFX allocations: counting and variants.
\newblock In \emph{Proceedings of AAAI}, pages 14036--14044, 2025.

\bibitem{Mahara2107}
R.~Mahara.
\newblock Extension of additive valuations to general valuations on the
existence of EFX.
\newblock arXiv:2107.09901, 2021.

\bibitem{ChaudhuryThree}
B.~R. Chaudhury, J.~Garg, and K.~Mehlhorn.
\newblock EFX exists for three agents.
\newblock \emph{Journal of the ACM}, 71(1), Article 4, 2024;
arXiv:2002.05119.

\bibitem{Lipton2004}
R.~J. Lipton, E.~Markakis, E.~Mossel, and A.~Saberi.
\newblock On approximately fair allocations of indivisible goods.
\newblock In \emph{Proceedings of EC}, 2004.

\bibitem{CGM2019}
B.~R. Chaudhury, T.~Kavitha, K.~Mehlhorn, and A.~Sgouritsa.
\newblock A little charity guarantees almost envy-freeness.
\newblock \emph{SIAM Journal on Computing}, 50(4):1336--1358, 2021.

\bibitem{Berger2022}
B.~Berger, A.~Cohen, M.~Feldman, and A.~Fiat.
\newblock (Almost full) EFX exists for four agents (and beyond).
\newblock In \emph{Proceedings of AAAI}, 2022.

\bibitem{Flener2002}
P.~Flener, A.~M. Frisch, B.~Hnich, Z.~Kiziltan, I.~Miguel, J.~Pearson, and
T.~Walsh.
\newblock Breaking row and column symmetries in matrix models.
\newblock In \emph{Proceedings of CP}, 2002.

\bibitem{Branzei2025}
S.~Br\^anzei.
\newblock Computing envy-free up to any good (EFX) allocations via local search.
\newblock arXiv:2510.05429, 2025.

\bibitem{Akrami2026SAT}
H.~Akrami, A.~Mayorov, K.~Mehlhorn, S.~Srinivas, and C.~Weidenbach.
\newblock A counterexample to EFX: $n\ge3$ agents, $m\ge n+5$ items,
submodular valuations via SAT-solving.
\newblock arXiv:2604.18216, 2026.

\bibitem{ORtools}
L.~Perron and F.~Didier.
\newblock CP-SAT, version 9.15.
\newblock Google, \url{https://developers.google.com/optimization/cp/cp_solver/},
2026.

\bibitem{Z3}
L.~de Moura and N.~Bj{\o}rner.
\newblock Z3: an efficient SMT solver.
\newblock In \emph{Proceedings of TACAS}, 2008.

\bibitem{PySAT}
A.~Ignatiev, A.~Morgado, and J.~Marques-Silva.
\newblock PySAT: a Python toolkit for prototyping with SAT oracles.
\newblock In \emph{Proceedings of SAT}, pages 428--437, 2018.

\bibitem{Sinz2005}
C.~Sinz.
\newblock Towards an optimal CNF encoding of Boolean cardinality constraints.
\newblock In \emph{Proceedings of CP}, pages 827--831, 2005.

\bibitem{CaDiCaL}
A.~Biere, T.~Faller, K.~Fazekas, M.~Fleury, N.~Froleyks, and F.~Pollitt.
\newblock CaDiCaL 2.0.
\newblock In \emph{Proceedings of CAV}, pages 133--152, 2024.

\bibitem{DRATtrim}
N.~Wetzler, M.~J. H. Heule, and W.~A. Hunt.
\newblock DRAT-trim: efficient checking and trimming using expressive clausal
proofs.
\newblock In \emph{Proceedings of SAT}, pages 422--429, 2014.

\end{thebibliography}

\end{document}
